\section{Ordinary backtracking on an exposed solution face}
\label{sec:face}
The zero-base model in Section~\ref{sec:zero} supplies a full exposed
solution face. The behavior of ordinary FW on such a face can be analyzed
under weaker regularity: a Lipschitz gradient on the entire polytope
suffices. The conclusion concerns the actual trajectory of the algorithm.
A corresponding pointwise assertion over arbitrary near-optimal points
is false. We write $\delta_k:=\delta(x_k)$ along a trajectory.

Throughout this section let $Q=\operatorname{conv}(V)$ for a finite
set of atoms, with diameter $D>0$. Let $f$ be convex and continuously
differentiable on an open convex neighborhood of $Q$. Assume
\begin{equation}\label{face:lip}
 \norm{\nabla f(x)-\nabla f(y)}\le L_{\rm sm}\norm{x-y}
 \qquad (x,y\in Q),\quad L_{\rm sm}>0,
\end{equation}
and
\begin{equation}\label{face:exposed}
 X^*=F=\arg\min_{x\in Q}\inner{g^*}{x},\qquad
 g^*=\nabla f(x^*)\quad(x^*\in X^*).
\end{equation}
The gradients at all minimizers coincide: for a fixed minimizer $x^*$,
convexity and constrained optimality imply
$\inner{\nabla f(x^*)}{y-x^*}=0$ for every $y\in X^*$; the
nonnegative convex function
$f(\cdot)-f(x^*)-\inner{\nabla f(x^*)}{\cdot-x^*}$ vanishes at
each such $y$, which is an interior point of the neighborhood, so its
gradient vanishes there. If $F=Q$, the initial gap is zero. Otherwise set
\begin{equation}\label{face:margin}
 m=\min_{v\in V\setminus F}\inner{g^*}{v-x^*}>0,
 \qquad n_F=|V\cap F|,
 \qquad \overline L_{\rm bt}=\max\{M,2L_{\rm sm}\}.
\end{equation}
The LMO is required to return an atom from this fixed finite set; no
active-set representation is needed by the algorithm. Weak sharpness
and finite identification of optimal linear-oracle outputs under this
geometry are classical~\cite{burkeferris}; the result below gives a
finite budget of gap-ratio defects for the inherited quadratic-
backtracking rule.

\begin{theorem}[Trajectory ratio and finite defect budget]
\label{face:ratio}
Run the ordinary backtracking method of Section~\ref{sec:model}, with an
exact atom LMO and the first-halving, rejection-doubling rule. Under
\eqref{face:lip}--\eqref{face:exposed}, every infinite nonterminal run
satisfies
\begin{equation}\label{face:ratiolimit}
 \frac{G_k}{\delta_k}\longrightarrow 1.
\end{equation}
Consequently $U_k/\delta_k\to1$ along each such run, since
$\delta_k\le U_k\le G_k$.
More quantitatively, for every $\eta\in(0,1]$, after the first source
with $\delta\le\tau_\eta$, where
\begin{equation}\label{face:threshold}
 \tau_\eta=\frac{\eta m^2}{2L_{\rm sm}D^2},\qquad
 P_\eta=n_F\ceil{\frac{L_{\rm sm}\overline L_{\rm bt}D^4}
                              {\eta m^2}},
\end{equation}
at most $P_\eta$ nonterminal updates have
$G>(1+\eta)\delta$. For $\eta=1/3$, each nonexceptional update that does not reach an
optimum accepts its first fractional trial and halves the inherited
curvature.
\end{theorem}
\begin{proof}
The $L_{\rm sm}$-smooth descent lemma holds on every feasible trial
segment, so a trial is accepted once its estimate reaches
$L_{\rm sm}$. Since the search starts at half the preceding accepted
estimate and doubles after rejection,
$L_k\le\max\{L_{k-1}/2,2L_{\rm sm}\}$. Induction from
$L_{-1}=M$ gives $L_k\le\overline L_{\rm bt}$ for every update.
Put
\[
 t(x)=\inner{g^*}{x-x^*},\qquad h(x)=\delta(x)-t(x),\qquad
 C=\frac{L_{\rm sm}D^2}{2m^2}.
\]
Both $h$ and its gradient vanish on $F$, and $h\ge0$. In any atom
representation $x=\sum_v\lambda_vv$, the mass $\theta$ outside $F$
satisfies $\theta\le t/m$. Normalizing the remaining face weights
gives a point in $F$ at distance at most $D\theta$ from $x$. If
$\theta=1$, the same bound follows from $t\ge m$ and
$\operatorname{dist}(x,F)\le D$. Thus
$\operatorname{dist}(x,F)\le Dt/m$. The descent lemma applied from a
nearest point in $F$, where $h=0$ and $\nabla h=0$, and
\eqref{face:lip} give
\begin{equation}\label{face:residual}
 0\le h(x)\le Ct^2,\qquad
 \norm{\nabla h(x)}\le\frac{L_{\rm sm}D}{m}t.
\end{equation}

Suppose the selected atom $s\in F$ has positive current weight
$\lambda_s$. For any $v\in V\cap F$, the transfer point
$z_v=x+\lambda_s(v-s)$ is feasible and has the same value of $t$.
Convexity and~\eqref{face:residual} imply
\[
 \lambda_s\inner{\nabla h(x)}{v-s}
 \le h(z_v)-h(x)\le Ct^2.
\]
Consequently the face contribution to
$\inner{\nabla h(x)}{x-s}$ is at most $Ct^2/\lambda_s$.
The nonface contribution is at most
$\theta D\norm{\nabla h(x)}\le 2Ct^2$. Since
$\inner{g^*}{x-s}=t$ and $\delta=t+h\ge t$, this yields
\begin{equation}\label{face:defect}
 0\le\frac{G-\delta}{\delta}
 \le 2Ct+\frac{C}{A_s},\qquad A_s=\frac{\lambda_s}{t},
\end{equation}
where the last term is interpreted as $+\infty$ if $\lambda_s=0$.
Here the lower bound $G\ge\delta$ is the standard convexity bound.

If $\delta\le\tau_\eta=\eta/(4C)$, then
$D\norm{\nabla f(x)-g^*}\le\eta m/2$. Every atom outside $F$ has a
strictly larger current gradient score than every atom in $F$; hence
all subsequent LMO outputs lie in $F$. Propagate any atom representation
by ordinary FW updates. A full update reaches an optimum. On each
remaining update, if $\alpha<1$ then
\[
 t_+=(1-\alpha)t,\qquad
 A_s^+=A_s+\frac{\alpha}{(1-\alpha)t},\qquad
 A_v^+=A_v\quad(v\ne s).
\]
Fractionalness and $G\ge\delta\ge t$ imply
\[
 A_s^+-A_s\ge\frac{\alpha}{t}
 =\frac{G}{L_k\norm{s-x}^2t}
 \ge\frac{1}{\overline L_{\rm bt}D^2}.
\]
Once $A_s\ge2C/\eta$, inequality~\eqref{face:defect} gives
$G\le(1+\eta)\delta$. Each face atom can be selected before reaching
this threshold at most
\[
 \ceil{\frac{2C\overline L_{\rm bt}D^2}{\eta}}
 =\ceil{\frac{L_{\rm sm}\overline L_{\rm bt}D^4}
                 {\eta m^2}}
\]
times. Summing over the $n_F$ face atoms proves the defect budget.

For each $\eta>0$, only finitely many later ratios exceed $1+\eta$.
To see that the run cannot stay away from the optimum, every accepted
update decreases $f$ by at least
$\min\{\delta_k/2,\delta_k^2/(2\overline L_{\rm bt}D^2)\}$,
using~\eqref{eq:acceptedprogress} for a full step and the fractional
step formula otherwise. The finite total decrease implies
$\delta_k\to0$. This proves~\eqref{face:ratiolimit}.

Finally let $\eta=1/3$ and consider a nonexceptional nonterminal
update after the stated entry. Its LMO atom is in $F$ and
$G\le4\delta/3$. If the first trial is fractional, convexity gives
$f(x+\alpha(s-x))\le f(x)-\alpha\delta$, while the quadratic
model requires a decrease of $\alpha G/2$; since $G<2\delta$, this
first fractional trial is accepted. If the first trial is full at
curvature $\ell$, acceptance is equivalent to
$\ell\norm{s-x}^2\ge2(G-\delta)$, because $f(s)=f^*$.
If it passes, the update terminates exactly. If it fails, the first
doubled curvature crossing this threshold is less than
$4(G-\delta)/\norm{s-x}^2\le G/\norm{s-x}^2$, where the last
inequality uses $G\le4\delta/3$. That doubled trial remains full and
is accepted, so it also terminates exactly. Hence every nonterminal
regular update accepts its first fractional trial, which halves the
inherited curvature.
\end{proof}

The normalized-weight cancellation in~\eqref{face:defect} is an
elementary rescaling; a related rescaling appears in
Harks~\cite[Section~3.2]{harks}. The budget controls the number of
exceptional updates, not the index of the last one.

\begin{corollary}[Observable local complexity]
\label{face:local}
Set
\[
 \tau=\frac{m^2}{6L_{\rm sm}D^2},\qquad
 P=n_F\ceil{\frac{3L_{\rm sm}\overline L_{\rm bt}D^4}{m^2}}.
\]
If $K$ is a source with $U_K\le\tau$, then after $N>P$ further
nonterminal updates,
\begin{equation}\label{face:certificate}
 U_{K+N}<\overline L_{\rm bt}D^2\,2^{-\ceil{(N-P)/(P+1)}}.
\end{equation}
In particular, for every $\varepsilon>0$, at most
\begin{equation}\label{face:localcount}
 P+(P+1)\max\left\{1,
 \ceil{\log_2\frac{\overline L_{\rm bt}D^2}{\varepsilon}}\right\}
\end{equation}
updates are needed after this entry until $U\le\varepsilon$. The entry
time is finite; a valid conservative bound is
\begin{equation}\label{face:entry}
 K\le\ceil{\frac{U_0}{\min\{\tau/2,\tau^2/(2\overline L_{\rm bt}D^2)\}}}.
\end{equation}
When the model assumptions of Theorem~\ref{thm:fw} are available, its
sharper phase bound can replace~\eqref{face:entry}.
\end{corollary}
\begin{proof}
There are at most $P$ exceptional updates, splitting the regular
updates into at most $P+1$ blocks. One block has length at least
$\ceil{(N-P)/(P+1)}$. A regular block of length $r$ ends with
curvature at most $\overline L_{\rm bt}2^{-r}$. At its final
fractional source, $U\le G< L_k\norm{s-x}^2\le
\overline L_{\rm bt}D^2 2^{-r}$. The certificate is nonincreasing
thereafter, proving~\eqref{face:certificate} and the stated local
count. A full update already solves the problem. Until entry,
$G\ge U>\tau$, so every accepted update decreases $f$ by at least
the denominator in~\eqref{face:entry}; the available initial decrease
is at most $U_0$. This proves the entry bound.
\end{proof}
The constants can be large, particularly when the face has many
atoms. The result gives an accuracy-independent finite defect budget
and logarithmic local certificate complexity, without a Hessian
modulus of continuity; it does not assert a small practical condition
number or a minimax bound. For $T$ updates and a final source check,
$N_g=N_{\rm LMO}=T+1$ and
\[
 N_f\le1+2T+\ceil{\log_2\max\{1,2L_{\rm sm}/M\}}.
\]

\subsection{A pointwise counterexample and an infinite exact trajectory}
\label{face:examples}
The pointwise extension of~\eqref{face:ratiolimit} to arbitrary feasible
points is false. On
\[
 Q=\operatorname{conv}\{(0,-1),(0,1),(1,0)\},\qquad
 f(t,z)=t+\tfrac12(at+z-1)_+^2,\quad a>1,
\]
the gradient is globally Lipschitz with constant $a^2+1$,
$X^*=F=\{0\}\times[-1,1]$, and $g^*=(1,0)$. At
$x_t=(t,1-t)$, the unique LMO atom is $(0,-1)$ and
\[
 \delta(x_t)=t+\tfrac12(a-1)^2t^2,\qquad
 G(x_t)=(2a-1)t+(a-1)^2t^2.
\]
Thus $G(x_t)/\delta(x_t)\to2a-1>1$. These arbitrary feasible
points need not be a tail of the algorithm's trajectory.

Conversely, ordinary FW need not terminate exactly even in the
zero-base model. Let
\begin{equation}\label{face:c2example}
 Q=[0,1]\times[-1,1],\qquad
 f(t,z)=t+c(t_+)^3(1+z^2/4),\qquad c=81/476.
\end{equation}
This function is $C^2$ and convex on an open rectangle containing
$Q$ with $|z|<\sqrt2$. For $t>0$ its Hessian determinant is
$3c^2t^4(1-z^2/2)$; for $t\le0$ its Hessian is zero. On $Q$ its
trace is at most $8ct\le8cf$, so it satisfies the original model
with $H_0=0$ and $H_1=8c$.
Start at $(t_0,1/3)$, where $0<t_0\le1$, with $M=9t_0/4$.
At every iterate the unique LMO selects the opposite endpoint of the
optimal face. Inductively,
\begin{equation}\label{face:c2trajectory}
 t_k=2^{-k}t_0,\qquad z_k=(-1)^k/3,\qquad
 L_k=9t_k/8,\qquad \alpha_k=1/2.
\end{equation}
Indeed, at either sign of $z=1/3$,
$G=t+(9/16)t^3$ and $\norm{s-x}^2=t^2+16/9$, so
$G/\norm{s-x}^2=9t/16$. The first half-curvature trial therefore
has step $1/2$. Since
$f(t,\pm1/3)=t+(333/1904)t^3$, the quadratic acceptance right
side minus the proposed endpoint value is
\[
 \frac{t}{4}+\frac{189}{15232}t^3>0.
\]
This proves~\eqref{face:c2trajectory} by induction. No iterate is
optimal and $\delta_{k+1}/\delta_k\to1/2$. The example is $C^2$;
no $C^\infty$ regularity is claimed.

\begin{corollary}[Sparse endpoint probes]
\label{face:endpoint}
Under the assumptions of Theorem~\ref{face:ratio}, add an
Armijo test
\begin{equation}\label{face:armijo}
 f(s)\le f(x)-\sigma G
\end{equation}
with any fixed $0<\sigma<1$ at a predetermined unbounded set of
iteration indices. A failed test leaves the current point and curvature memory unchanged
before the ordinary backtracking fallback. An accepted test sets the
next point to $s$ and carries forward the current positive, finite
curvature memory. With exact oracles and arithmetic, this algorithm
reaches an exact optimum after finitely many updates.
\end{corollary}
\begin{proof}
Every accepted nonoptimal endpoint strictly decreases the objective
to one of the finitely many atom values, so there can be only finitely
many such acceptances. After the last one, an infinite run consists
of the ordinary fallback steps. Eventually $s\in X^*$ and
$G/\delta\to1$, hence $\delta\ge\sigma G$; the next scheduled
endpoint test succeeds at an optimum. If the run terminates earlier,
the claim already holds.
\end{proof}
This is an exact-arithmetic termination result with no positive
stopping tolerance or imposed iteration cap. It is not a speedup
claim. A persistent cache of exact atom values limits the additional
cost of wholly rejected probes to at most $|V|$ over any fixed number
of fallback updates; the empirical qualification is recorded in
Appendix~\ref{app:endpoint}.
